%======================================================================
% فصل سوم: طراحی سیستم
%======================================================================
\chapter{طراحی سیستم}
\label{ch:design}

\section{مقدمه}

پس از بررسی مسئله و پیشینه، در این فصل طراحی سامانه‌ی پیاده‌سازی‌شده
تشریح می‌شود. هدف از این طراحی، ساخت بستری است که بتواند هم به عنوان
نقطه‌ی ورود یکپارچه برای درخواست‌های هوش مصنوعی عمل کند و هم قابلیت‌های
مدیریتی، مالی و نظارتی لازم برای بهره‌برداری سازمانی را فراهم آورد.

در طراحی سامانه، سه اصل راهنما در نظر گرفته شده است: نخست،
\textbf{سازگاری بیرونی} با رابط \lr{OpenAI} برای کاهش هزینه‌ی مهاجرت
مشتریان؛ دوم، \textbf{تفکیک مسئولیت‌ها} در لایه‌های داخلی برای افزایش
نگهداشت‌پذیری؛ و سوم، \textbf{قابلیت کنترل و گزارش‌گیری} برای مدیریت
مصرف، هزینه و امنیت.

\section{نیازمندی‌های طراحی}

پیش از تشریح معماری، لازم است نیازمندی‌هایی که طراحی باید به آن‌ها پاسخ
دهد به‌صورت صریح فهرست شوند. جدول \ref{tab:design-requirements} این
نیازمندی‌ها را در دو دسته‌ی کارکردی و غیرکارکردی خلاصه می‌کند.

\begin{table}[H]
\centering
\caption{نیازمندی‌های اصلی مؤثر بر طراحی سامانه}
\label{tab:design-requirements}
\renewcommand{\arraystretch}{1.4}
\begin{tabularx}{\textwidth}{>{\raggedleft\arraybackslash}p{2.2cm} >{\raggedleft\arraybackslash}p{2.8cm} X}
\toprule
\textbf{شناسه} & \textbf{نوع} & \textbf{شرح} \\
\midrule
\lr{F1} & کارکردی & ارائه‌ی نقاط پایانی سازگار با \lr{OpenAI} برای گفتگو، \lr{responses}، جاسازی متن، تصویر و صوت \\
\lr{F2} & کارکردی & پشتیبانی از چند ارائه‌دهنده با امکان نگاشت تفاوت‌های قراردادی در لایه‌ی درونی \\
\lr{F3} & کارکردی & احراز هویت کاربر، مدیریت مشتریان \lr{API} و کلیدهای برنامه‌ای \\
\lr{F4} & کارکردی & کنترل مالی شامل کیف پول، اشتراک، فاکتور و ثبت دقیق مصرف \\
\lr{F5} & کارکردی & گزارش‌گیری از مصرف، تراکنش‌ها و عملیات حساس \\
\lr{NF1} & غیرکارکردی & کاهش وابستگی مشتری به ارائه‌دهنده و تسهیل مهاجرت میان سرویس‌ها \\
\lr{NF2} & غیرکارکردی & تفکیک مسئولیت، آزمون‌پذیری و قابلیت توسعه‌ی تدریجی \\
\lr{NF3} & غیرکارکردی & کنترل امنیتی شامل محدودسازی نرخ، فهرست سفید \lr{IP} و ثبت ممیزی \\
\lr{NF4} & غیرکارکردی & خارج کردن کارهای غیر بحرانی از مسیر اصلی درخواست با تکیه بر کش و صف \\
\bottomrule
\end{tabularx}
\end{table}

چند نیازمندی نیز آگاهانه از دامنه‌ی نسخه‌ی فعلی خارج شده‌اند. برای نمونه،
انتخاب هوشمند ارائه‌دهنده، اتصال \lr{OTP} به سرویس واقعی پیامک یا ایمیل و
یکپارچه‌سازی با درگاه پرداخت بیرونی به‌عنوان جهت توسعه‌ی آتی نگه داشته
شده‌اند. بنابراین طراحی فعلی باید به‌گونه‌ای باشد که این قابلیت‌ها را در
آینده بدون بازنویسی گسترده بپذیرد.

\section{معماری کلی سیستم}

سامانه از چند زیرسامانه‌ی مکمل تشکیل شده است که هر کدام مسئول بخش
مشخصی از چرخه‌ی خدمت‌رسانی هستند. نمای منطقی این اجزا در
جدول \ref{tab:main-components} خلاصه شده است.

\begin{table}[H]
\centering
\caption{اجزای اصلی معماری سامانه}
\label{tab:main-components}
\renewcommand{\arraystretch}{1.4}
\begin{tabularx}{\textwidth}{>{\raggedleft\arraybackslash}p{3.5cm} X}
\toprule
\textbf{جزء} & \textbf{وظیفه} \\
\midrule
\lr{REST API Layer} & ارائه‌ی نقاط پایانی نسخه‌بندی‌شده برای احراز هویت، مدیریت کاربران، کلیدهای \lr{API}، کیف پول، اشتراک، گزارش و درخواست‌های هوش مصنوعی \\
\lr{Gateway Core} & اعتبارسنجی، انتخاب ارائه‌دهنده و مدل، تخمین مصرف، فراخوانی ارائه‌دهنده، یکسان‌سازی پاسخ و ثبت مصرف \\
\lr{Billing Core} & مدیریت کیف پول، تراکنش‌ها، اشتراک‌ها، فاکتورها و اعمال هزینه‌ی هر درخواست \\
\lr{Identity \& Access} & احراز هویت مبتنی بر \lr{OTP}، مدیریت نقش‌ها، مشتریان \lr{API} و کلیدها \\
\lr{Reporting} & تولید گزارش مصرف، دفترکل کیف پول، گزارش فاکتورها و انباره‌سازی تجمیعی \\
\lr{SQL Server} & نگهداری داده‌های پایدار شامل کاربران، مدل‌ها، مصرف، تراکنش‌ها و گزارش‌ها \\
\lr{Redis} & کش داده‌های پرتکرار، شمارنده‌های محدودسازی نرخ و صف کارهای پس‌زمینه \\
\lr{Queue Workers} & اجرای غیرهمگام ارسال \lr{OTP}، تولید \lr{PDF} فاکتور، تجمیع روزانه و ارسال گزارش خارجی \\
\lr{Playground UI} & محیط آزمایشی برای ارسال تعاملی درخواست و مشاهده‌ی پاسخ سامانه \\
\bottomrule
\end{tabularx}
\end{table}

از دید تعامل، سامانه سه مسیر اصلی دارد. در مسیر مدیریتی، کاربر انسانی پس
از احراز هویت با \lr{OTP} به داشبورد و نقاط پایانی مدیریتی دسترسی پیدا
می‌کند. در مسیر برنامه‌ای، نرم‌افزارهای مشتری با کلید \lr{API} درخواست‌های
هوش مصنوعی را به نقاط پایانی سازگار با \lr{OpenAI} می‌فرستند. مسیر سوم نیز
به کارهای پس‌زمینه اختصاص دارد؛ کارهایی مانند ارسال گزارش خارجی و تولید
\lr{PDF} که زمان‌بر یا غیر بحرانی هستند از طریق صف اجرا می‌شوند.

\subsection{جریان درخواست}

جریان منطقی یک درخواست هوش مصنوعی در سامانه از دریافت آن در نقطه‌ی پایانی
\lr{/api/v1/ai/...} آغاز می‌شود. سپس کلید \lr{API} و وضعیت کاربر مرتبط
استخراج و اعتبارسنجی می‌شود، محدودیت نرخ و مجاز بودن آدرس \lr{IP} بررسی
می‌گردد و ارائه‌دهنده و مدل بر اساس مقادیر صریح \code{provider} و
\code{model} تعیین می‌شوند. پس از آن، سامانه پیش از تماس با
ارائه‌دهنده مصرف و هزینه‌ی احتمالی را تخمین می‌زند و وضعیت اشتراک یا
موجودی کیف پول را کنترل می‌کند. اگر درخواست مجاز باشد، فراخوانی از طریق
آداپتور مناسب انجام می‌شود، پاسخ یکدست‌سازی می‌گردد، مصرف واقعی از آن
استخراج یا با تخمین تلفیق می‌شود، رکوردهای گزارش و مصرف در پایگاه داده
ثبت می‌شوند و در پایان هزینه اعمال و در صورت نیاز گزارش خارجی به‌صورت
غیرهمگام ارسال می‌شود.

این جریان به‌گونه‌ای طراحی شده است که خطاهای مالی یا امنیتی پیش از ارسال
درخواست به ارائه‌دهنده تشخیص داده شوند و از هزینه‌ی غیرضروری جلوگیری شود.

\section{مدل داده}

مدل داده‌ی سامانه بر اساس چهار حوزه‌ی اصلی طراحی شده است: هویت و دسترسی،
ارائه‌دهندگان و مدل‌ها، مالی و صورت‌حساب، و گزارش‌گیری. این تقسیم‌بندی با
تفکیک ماژول‌های نرم‌افزار نیز هم‌راستا است.

\subsection{جداول احراز هویت و دسترسی}

هسته‌ی مدیریت کاربران بر جدول \lr{users} استوار است و با جداول
\lr{roles}، \lr{permissions} و جداول واسط مربوطه تکمیل می‌شود. برای ورود
کاربر به سامانه، جدول \lr{otp\_challenges} نگهدارنده‌ی مقصد، کانال، کد
هش‌شده، زمان انقضا و تعداد دفعات تلاش است.

برای دسترسی برنامه‌ای نیز دو سطح مجزا در نظر گرفته شده است. سطح نخست
\lr{api\_clients} است که یک برنامه یا مصرف‌کننده‌ی منطقی متعلق به کاربر
را نمایش می‌دهد. سطح دوم \lr{api\_keys} است که پیشوند کلید، هش کلید،
حوزه‌های دسترسی، محدودیت نرخ، فهرست \lr{IP}‌های مجاز و تاریخ انقضا یا
ابطال را نگهداری می‌کند.

این طراحی باعث می‌شود کاربر بتواند برای چند نرم‌افزار مختلف، کلیدهای
مجزا با سیاست‌های امنیتی متفاوت ایجاد کند.

\subsection{جداول ارائه‌دهندگان و مدل‌ها}

پیکربندی ارائه‌دهندگان در جدول \lr{providers} و مشخصات مدل‌ها در جدول
\lr{provider\_models} ذخیره می‌شود. هر ارائه‌دهنده شامل نام، نوع،
\lr{Base URL}، وضعیت، اولویت و پیکربندی محرمانه است. هر مدل نیز با
\code{model\_key}، قابلیت‌ها و پیکربندی قیمت‌گذاری توصیف می‌شود.

در کنار این دو جدول، جدول \lr{routing\_rules} نیز برای توسعه‌ی آینده در
نظر گرفته شده است. هرچند در نسخه‌ی فعلی، انتخاب ارائه‌دهنده و مدل به
صورت صریح توسط مشتری انجام می‌شود، وجود این جدول امکان پیاده‌سازی
راهبردهای پیشرفته‌تر را در آینده فراهم می‌کند.

\subsection{جداول صورت‌حساب}

بخش مالی سامانه از چند جدول مکمل تشکیل شده است. جدول \lr{wallets} موجودی
اصلی هر کاربر و واحد پول را نگه می‌دارد، \lr{wallet\_transactions} دفترکل
تراکنش‌های بدهکار و بستانکار است، \lr{subscription\_plans} طرح‌های اشتراک
و قیمت و امکانات آن‌ها را تعریف می‌کند، \lr{subscriptions} وضعیت اشتراک
فعال کاربر و بازه‌ی اعتبار آن را ثبت می‌نماید و در نهایت جداول
\lr{invoices} و \lr{invoice\_items} برای نگهداری فاکتورها و اقلام هر
فاکتور به‌کار می‌روند.

این ساختار دو الگوی مالی را به‌صورت هم‌زمان پشتیبانی می‌کند: مصرف بر اساس
کیف پول و مصرف مبتنی بر اشتراک.

\subsection{جداول گزارش مصرف}

هر درخواست عبوری از سامانه در صورت فعال بودن ثبت گزارش، در جدول
\lr{gateway\_requests} ذخیره می‌شود. سپس اقلام مصرف قابل‌صورتحساب در جدول
\lr{usage\_records} ثبت می‌گردند. برای کاهش هزینه‌ی گزارش‌گیری بلندمدت،
خلاصه‌های روزانه در جدول \lr{daily\_usage\_rollups} تولید می‌شود.

دو جدول مکمل دیگر نیز در این حوزه وجود دارند: \lr{audit\_logs} برای ثبت
عملیات حساس، مانند ایجاد یا ابطال کلید، و \lr{provider\_health\_checks}
برای پایش سلامت ارائه‌دهندگان.

\section{طراحی \lr{API}}

طراحی \lr{API} بر دو اصل استوار است: نسخه‌بندی شفاف و تفکیک نقاط پایانی
مدیریتی از نقاط پایانی هوش مصنوعی. همه‌ی مسیرها زیر پیشوند \lr{/api/v1}
قرار گرفته‌اند تا توسعه‌ی نسخه‌های بعدی بدون شکستن سازگاری ممکن باشد.

\subsection{نقاط پایانی احراز هویت و مدیریتی}

نقاط پایانی مدیریتی از توکن نشست \lr{Sanctum} استفاده می‌کنند
~\cite{sanctum2024}. این مسیرها در عمل چند گروه را پوشش می‌دهند: مسیرهای
احراز هویت مانند \lr{/auth/otp/start} و \lr{/auth/otp/verify}، مسیر
پروفایل کاربر در \lr{/me}، مسیرهای مدیریت مشتری و کلید در
\lr{/api-clients} و \lr{/api-keys/...}، مسیرهای مالی مانند
\lr{/wallet}، \lr{/plans}، \lr{/subscriptions} و \lr{/invoices}،
مسیرهای گزارش‌گیری مانند \lr{/reports/usage}،
\lr{/reports/wallet-ledger} و \lr{/reports/invoices}، و در نهایت مسیرهای
مدیریت ارائه‌دهندگان در \lr{/providers} و
\lr{/providers/\{provider\}/models}.

علاوه بر این، دریافت \lr{PDF} فاکتور با مسیر \lr{/invoices/\{invoice\}/pdf}
و از طریق \term{پیوند امضاشده}{Signed URL} طراحی شده است تا دسترسی کنترل‌شده
و موقت به فایل فراهم شود.

\subsection{نقاط پایانی هوش مصنوعی}

برای بیشینه کردن سازگاری با مشتریان موجود، شش نقطه‌ی پایانی اصلی مطابق
قرارداد بیرونی \lr{OpenAI} در نظر گرفته شده است: \lr{/ai/responses}،
\lr{/ai/chat/completions}، \lr{/ai/embeddings}،
\lr{/ai/images/generations}، \lr{/ai/audio/transcriptions} و
\lr{/ai/audio/speech}.

همه‌ی این درخواست‌ها با کلید \lr{API} احراز هویت می‌شوند و به‌صورت صریح
دارای دو فیلد \code{provider} و \code{model} هستند. این تصمیم طراحی دو
مزیت دارد: نخست، از ابهام در قیمت‌گذاری جلوگیری می‌کند؛ دوم، مسئولیت
انتخاب ارائه‌دهنده را برای مشتری شفاف می‌سازد.

\section{الگوی خط لوله}

هسته‌ی طراحی سامانه بر \term{الگوی خط لوله}{Pipeline Pattern} استوار است.
در این الگو، هر مرحله یک مسئولیت مشخص دارد و ورودی و خروجی آن در قالب
یک \term{زمینه‌ی درخواست}{Request Context} مشترک حمل می‌شود. مزایای اصلی
این الگو در پروژه‌ی حاضر روشن است: جریان پردازش خواناتر می‌شود و
نگرانی‌های مختلف از هم جدا می‌مانند، هر مرحله را می‌توان به‌صورت مستقل
آزمود، افزودن یا حذف مرحله‌ها بدون دستکاری گسترده‌ی کد ممکن می‌شود و
مدیریت خطا و توقف زودهنگام در صورت بروز مشکل ساده‌تر خواهد بود.

در نسخه‌ی فعلی، خط لوله‌ی درگاه شامل چهارده مرحله است:

\begin{enumerate}
  \item \code{ResolveApiKeyPipe}
  \item \code{EnforceIpAllowlistPipe}
  \item \code{ValidateGatewayPayloadPipe}
  \item \code{RateLimitPipe}
  \item \code{SelectProviderPipe}
  \item \code{ValidateProviderSelectionPipe}
  \item \code{EstimateUsagePipe}
  \item \code{CheckSubscriptionOrWalletPipe}
  \item \code{DispatchProviderRequestPipe}
  \item \code{NormalizeProviderResponsePipe}
  \item \code{MeterUsagePipe}
  \item \code{PersistLogsPipe}
  \item \code{ChargeUsagePipe}
  \item \code{DispatchExternalLogsPipe}
\end{enumerate}

ترتیب این مراحل عمدی است. برای نمونه، تخمین مصرف باید پیش از بررسی کیف
پول انجام شود، و ثبت گزارش نیز باید پس از استخراج مصرف واقعی انجام گیرد.

\section{استراتژی کش}

در طراحی سامانه، کش با راهبرد \term{کش‌کنار}{Cache-Aside} در نظر گرفته شده
است. در این الگو، داده ابتدا از کش خوانده می‌شود و در صورت نبود، از
پایگاه داده بازیابی و سپس در کش ذخیره می‌گردد. این راهبرد برای داده‌های
کم‌تغییر و پرتکرار مناسب است.

مهم‌ترین اقلامی که در طراحی سامانه برای کش در نظر گرفته شده‌اند شامل فهرست
طرح‌های فعال با کلید \code{plans:active}، فهرست ارائه‌دهندگان با کلید
\code{providers:all}، پیکربندی هر ارائه‌دهنده با الگوی
\code{providers:config:providerName}، مدل‌های هر ارائه‌دهنده با الگوی
\code{provider_models:list:providerId} و پروفایل کاربر با الگوی
\code{user:profile:userId} هستند.

زمان اعتبار هر گروه از طریق پیکربندی مرکزی تعیین می‌شود تا بتوان سیاست
کش را بدون تغییر در منطق سرویس‌ها تنظیم کرد. استفاده از \lr{Redis} به
عنوان بستر کش، علاوه بر سرعت مناسب، امکان هم‌زیستی طبیعی کش، محدودسازی
نرخ و صف را نیز فراهم می‌کند~\cite{redis2024}.

\section{طراحی سیستم صورت‌حساب}

صورت‌حساب‌گذاری در سامانه بر پایه‌ی مدل‌های قیمت‌گذاری ذخیره‌شده در
فیلد \lr{pricing\_config} از جدول \lr{provider\_models} طراحی شده است.
این تصمیم باعث می‌شود که منطق قیمت، به‌جای پراکندگی در کد، به‌صورت
داده‌محور و قابل تغییر مدیریت شود.

برای هر درخواست، ابتدا یک تخمین از مصرف انجام می‌شود تا امکان پیش‌بررسی
موجودی وجود داشته باشد. سپس پس از دریافت پاسخ، مصرف واقعی از پاسخ
ارائه‌دهنده استخراج یا با تخمین پیشین تلفیق می‌شود. هزینه‌ی نهایی هر
درخواست از رابطه‌ی زیر به دست می‌آید:

\[
C = \sum_{i=1}^{n} q_i \times p_i
\]

که در آن \(q_i\) مقدار هر معیار مصرف و \(p_i\) هزینه‌ی واحد همان معیار است.
این معیارها می‌توانند شامل توکن ورودی، توکن خروجی، تعداد تصویر، ثانیه‌ی
صوت و تعداد نویسه‌های ورودی برای \lr{TTS} باشند.

در طراحی مالی، دو مسیر هم‌زمان پشتیبانی می‌شوند. اگر کاربر اشتراک فعال با
اعتبار مشمول مصرف داشته باشد، درخواست بدون برداشت مستقیم از کیف پول عبور
می‌کند. در غیر این صورت، هزینه از موجودی کیف پول کسر می‌شود و تراکنش
متناظر در دفترکل ثبت می‌گردد.

فاکتورها نیز به‌عنوان نمایش رسمی اقلام هزینه طراحی شده‌اند و اقلام آن‌ها
می‌توانند هم مصرف جاری و هم هزینه‌های اشتراک را پوشش دهند. این طراحی،
امکان گزارش‌گیری مالی و تولید \lr{PDF} رسمی برای هر دوره را فراهم می‌کند.
